
    \documentclass[a4paper]{article}
    \usepackage{savetrees}
    \usepackage{hyperref}
    \usepackage{pdfpages}

    \begin{document}

    \underline{\href{https://nbalepur.github.io/}{Nishant Balepur}}, Bhavya Rajasekaran, Hyunjin Jane Oh, Michael Xie, Atrey Desai, Vipul Gupta, Steven James Moore, Eunsol Choi, Rachel Rudinger, and {\bf Jordan Boyd-Graber}.  {\bf \href{http://cs.umd.edu/~jbg//docs/2026_acl_benchmarker.pdf}{BenchMarker: An Education-Inspired Toolkit for Highlighting Flaws in Multiple-Choice Benchmarks}}.  \emph{Association for Computational Linguistics}, 2026.



    \vspace{5cm}

    \begin{verbatim}
    @inproceedings{Balepur:Rajasekaran:Oh:Xie:Desai:Gupta:Moore:Choi:Rudinger:Boyd-Graber-2026,
	Title = {BenchMarker: An Education-Inspired Toolkit for Highlighting Flaws in Multiple-Choice Benchmarks},
	Author = {Nishant Balepur and Bhavya Rajasekaran and Hyunjin Jane Oh and Michael Xie and Atrey Desai and Vipul Gupta and Steven James Moore and Eunsol Choi and Rachel Rudinger and Jordan Lee Boyd-Graber},
	Booktitle = {Association for Computational Linguistics},
	Year = {2026},
	Url = {http://cs.umd.edu/~jbg//docs/2026_acl_benchmarker.pdf},
}

    \end{verbatim}

    {\bf Accessible Abstract:} Multiple-choice questions are a standard way to evaluate NLP systems, but they are riddled with flaws that limit their validity. Extending our previous position paper, we draw on educational testing theory to design BenchMarker, a toolkit that detects faulty MCQs that exist on the Internet, have guessable shortcuts, and writing issues that confuse students and LLMs. We show how BenchMarker can detect and help fix flaws in NLP benchmarks.\\~

    

    \noindent Links:
\begin{itemize}
\item \href{https://github.com/nbalepur/BenchMarker}{Code+Data} [\url{https://github.com/nbalepur/BenchMarker}]

\end{itemize}\

    \vspace{3cm}
    Downloaded from \url{http://cs.umd.edu/~jbg/docs/2026_acl_benchmarker.pdf}  \\

    \vspace{2cm}
    \textit{Contact Jordan Boyd-Graber (jbg@boydgraber.org) for questions about this paper.}

    \includepdf[pages={-}]{src_docs/2026_acl_benchmarker.pdf}

    \end{document}
    